\documentclass[11pt,a4paper]{exam}
\usepackage[utf8]{inputenc}
\usepackage{amsmath, amsthm, amssymb, amsfonts} %ams packages
\usepackage{graphicx, color} %grphics
\usepackage{enumitem} %personalized lists 
\usepackage[top=4cm, bottom=2cm, left=2cm, right=2cm, headsep = 3pt, footskip=2pt]{geometry} %pagelayout
\usepackage{multirow,multicol} %to use multirow on tables
\usepackage{tabularx}
\usepackage{array} %bmatrix and similar stuff
% \usepackage{hyperref}
\usepackage[slovene]{babel}

\usepackage{multicol}


%%%%%%%%%%%%%%%%%%%%%%%%%%%%%%%%%%%%%%%%%%%%%%%%%%%%%%%%%%%%%%%%%%%%%
%personalized commands to make latex look better (optional)
%%%%%%%%%%%%%%%%%%%%%%%%%%%%%%%%%%%%%%%%%%%%%%%%%%%%%%%%%%%%%%%%%%%%%%%%%%%
\renewcommand{\leq}{\leqslant} %menores y mayores iguales decentes
\renewcommand{\geq}{\geqslant}
\renewcommand{\subset}{\subseteq} 
\renewcommand{\supset}{\supseteq} 
\renewcommand{\subsetneq}{\varsubsetneq}
\renewcommand{\{}{\lbrace}
\renewcommand{\}}{\rbrace}
\newcommand{\sm}{\setminus} %setminus corto

%\renewcommand{\bar}{\overline}
%\renewcommand{\hat}{\widehat}

%%%%%%%%%%%%%%%%%%%%%%%%%%%%%%%%%%%%%%%%%%%%%%%%%%%%%%%%%%%%%%%%%%%%%}

%%%%%%%%%%%%%%%%%%%%%%%%%%%%%%%%%%%%%%%%%%%%%%%%%%%%%%%%%%%%%%%%%%%%%%%%%%%
%exam class commands
%%%%%%%%%%%%%%%%%%%%%%%%%%%%%%%%%%%%%%%%%%%%%%%%%%%%%%%%%%%%%%%%%%%%%%%%%%%
\pagestyle{head}
% \chead{ Matematična analiza 2024/2025 \\ Univerza v Ljublani, Pedagoška fakulteta \\ Vaje 1 \\17. februar 2025} %self-explanatory
% \runningheader{}{ Abstraktna algebra 2024/2025 \\ Univerza v Ljublani, Pedagoška fakulteta \\2. Kolokvij \\20. januar 2025}{} %self-explanatory
% \runningheader{}{}{}
% \footer{}{}{}
% \firstpagefooter{MATH 1190 3.0 M  }{Page \thepage\ of \numpages}{Copyright 2020 \textcopyright Antonio Montero}
% \firstpagefooter{}{}{}
% \runningfooter{MATH 1190 3.0 M }{Page \thepage\ of \numpages}{Copyright 2020 \textcopyright Antonio Montero}
% \firstpagefootrule 
% \runningfootrule

%%Solutions
\unframedsolutions
% \printanswers
\SolutionEmphasis{\itshape}


%Personalized commands
\pointpoints{točka}{točk}
% \gradetable[h][questions]
\addpoints
\hqword{Naloga}
% \vpgword{text}
\hpword{Točke}
% \hsword{}
% \hbpword{}
\htword{Skupaj}

% \pointsinrightmargin
% \extrawidth{-2cm}
\marginpointname{ \points}


\chqword{Naloga}
\chpgword{Stran}
\chpword{točke}
% \chbpword{Puntos extra}
% \chsword{Puntos obtenidos
\chtword{Skupaj}
%%%%%%%%%%%%%%%%%%%%%%%%%%%%%%%%%%%%%%%%%%%%%%%%%%%%%%%%%%%%%%%%%%%%%

% \graphicspath{{../img/}} %Uso: \graphicspath{{RelativePath}}

%%%%%%%%%%%%%%%%%%%%%%%%%%%%%%%%%%%%%%%%%%%%%%%%%%%%%%%%%%%%%%%%%%%%%
%Personalized commands

%Personalized commands

\usepackage{tabularx}
\newcolumntype{C}{>{\centering\arraybackslash}X}%
\newcolumntype{R}{>{\raggedleft\arraybackslash}X}%
\newcolumntype{L}{>{\raggedright\arraybackslash}X}%

\newcolumntype{\cc}[1]{>{\hsize=#1\hsize\raggedright\arraybackslash}X}%
\newcolumntype{\rr}[1]{>{\hsize=#1\hsize\raggedleft\arraybackslash}X}%
\newcolumntype{\ll}[1]{>{\hsize=#1\hsize\centering\arraybackslash}X}%

\newcommand{\TT}{\mathbf{T}}
\newcommand{\FF}{\mathbf{F}}

\DeclareMathOperator{\sh}{sh}
\DeclareMathOperator{\ch}{ch}
\let\th\relax
\DeclareMathOperator{\th}{th}
\DeclareMathOperator{\arsh}{arsh}
\DeclareMathOperator{\arch}{arch}
\DeclareMathOperator{\arth}{arth}
\DeclareMathOperator{\arctg}{arctg}

%%%%%%%%%%%%%%%%%%%%%%%%%%%%%%%%%%%%%%%%%%%%%%%%%%%%%%%%%%%%%%%%%%%%%
		\newif\ifmarking
% 		\markingtrue %comment to hide marking suggestions

		\newenvironment{marking}
		{\hrulefill
		 
		\textbf{Marking suggestions:} 
		\ttfamily }
		{\hrulefill }

\allowdisplaybreaks

% \renewcommand{\thepartno}{\roman{partno}}
\renewcommand{\partlabel}{(\textit{\thepartno})}
%%%%%%%%%%%%%%%%%%%%%%%%%%%%%%%%%%%%%%%%%%%%%%%%%%%%%%%%%%%%%%%%%%%%%

\begin{document}
\chead{Matematična analiza 2024/2025 \\ Univerza v Ljublani, Pedagoška fakulteta \\ 1. Kolokvij \\3. april 2025}
% \footer{}{}{}
\pagestyle{head}
\begin{questions}
	\question
	 Dana je funkcija $f(x)=(2-x^2)e^{-x}$.
	 \begin{parts}
		\part Poiščite ničle in pole funkcije $f$ (če obstajajo), zapišite njeno definicijsko območje in raziščite obnašanje funkcije na robovih območja.
		\part Poiščite intervale naraščanja in padanja ter stacionarne točke funkcije $f$. Za stacionarne točke  določite, ali so lokalni/globalni minimum, maksimum ali nobena od njih.
		\part Poiščite intervale konveksnosti in konkavnosti ter prevoje funkcije $f$.
		\part Določite, ali ima funkcija asimptote. Če je tako, jih določite.
		\part Čimbolj natančno skicirajte graf funkcije $f$ s pomočjo prejšnjih točk.
	 \end{parts}
	
	 \vspace{1 cm}

	%  \question 
	%  \begin{parts}	
	% 	\part Pod kakšnim kotom se sekata grafa funkcij $f(x) = \sin(x)$ in $g(x)=\cos(2x)$?
	% 	\part Pokažite, da se krivulji $x^{2} + y^{2} = a$ in $xy=b$ sekata pravokotno za vsak $a,b \in \mathbb{R}$.
	% \end{parts}

	\question Dani sta funkciji $f(x)=\arctan (2x)$ and $g(x)=\frac{1}{2} \arctan \left( \frac{4x}{1-4x^{2}} \right)$ za $x \in (1,\infty)$. 
	\begin{parts}
		\part Dokažite, da imata $f$ in $g$ enaka odvoda povsod na intervalu $(1,\infty)$. 
		\part Ali sta tudi $f$ in $g$ enaki? Če ne, ugotovite kakšna zveza velja   med njima.
	\end{parts}
		
		\vspace{1cm}
	\question Iz kvadratnega kartona velikosti $30$ cm bomo izdelali škatlo (brez pokrova) tako, da bomo iz vogalov izrezali kvadrate in prepognili robove dobljenega kosa v obliki križa. Poišči mere take škatle z največjo prostornino.
	
	\vspace{1cm}
	\question Izračunajte naslednje integrale
	\begin{parts}
		\part $\displaystyle \int \sin^{5}(x) dx$
		\part $\displaystyle \int x \ln(1 + x^2) dx$

		

	\end{parts}
\end{questions}




\end{document}
